
\def \Subject {حملات خصمانه}
\def \Session { دو}
% \setcounter{chapter}{\Session-1}
\renewcommand{\chaptername}{سوال}
\chapter{\Subject}

\section{ مرور مفاهیم}

\subsection{دسته‌بندی روش‌های حمله متخاصم}
\begin{itemize}
	\item 
\textbf{جعبه سیاه:}
در حمله جعبه سیاه، مهاجم اطلاعات خاصی از مدل ندارد و فقط می‌تواند به مدل ورودی بدهد و برچسب و یا احتمال کلاس‌ها بگیرد. 
\item 
\textbf{جعبه سفید:}
در حمله جعبه سفید مهاجم اطلاعات مدل را دارد (معماری، وزن‌ها و غیره) و می‌تواند با استفاده از این اطلاعات سعی کند مدل را به اشتباه بیاندازد. 
\end{itemize}

یک دسته‌بندی که به نوعی میان این دو است حمله خاکستری است که در آن مهاجم اطلاعات محدودی از مدل را دارد. برای مثال معماری مدل را می‌داند ولی وزن‌های آن را خیر.

\subsection{اهمیت حملات جعبه سیاه}
در حملات جعبه سیاه، مهاجم نیاز به هیچ اطلاعاتی از مدل ندارد و در نتیجه در هر حالتی قابل انجام است. اگر یک مدل نسبت به حملات جعبه سفید آسیب‌پذیر باشد، ممکن است با بسته کردن مدل و حفاظت از اطلاعات مربوط به آن، برای مدل مشکلی پیش نیاید. ولی اگر مدلی نسبت به حمله جعبه سیاه آسیب‌پذیر باشد، صرف اینکه مهاجم بتواند به مدل ورودی بدهد می‌تواند منجر به بروز مشکل شود.
\subsection{حملات هدف‌دار و غیرهدف‌دار}
\begin{itemize}
	\item 
	\textbf{ غیرهدف‌دار:}
در این نوع حملات، تمرکز مهاجم صرفا بر این است که ورودی‌ای بدهد که مدل را به اشتباه بیاندازد.
	\item 
	\textbf{هدف دار :}
در این حملات، تمرکز مهاجم بر این است که ورودی‌ای بدهد که  مدل را به تولید خروجی مشخص و مورد نظر مهاجم وادار کند.
\end{itemize}
\section{حمله جعبه سیاه}
\subsection{آسیب‌پذیری نسبت به حمله جعبه سفید}
در این مقاله بیان می‌شود که طبق مقاله 
\cite{athalye2018obfuscated}
تقریبا همه مدل‌های دسته‌بندی تصویر نسبت به حملات جعبه سفید آسیب‌پذیرند. و تفسیر نویسنده مقاله این است که تصاویر طبیعی به مرزهای تصمیم مدل‌ها نزدیک هستند و این آسیب‌پذیری در ابعاد بالا تقریبا اجتناب‌ناپذیر است.
\subsection{روش حمله \lr{SimBA}}
همانطور که در بخش قبل بیان شد، تصاویر طبیعی به مرزهای تصمیم نزدیک هستند. همچنین اگر در نظر بگیریم که یک نمونه وجود دارد که با تغییراتی که برای انسان قابل تشخیص نیست، باعث شود کلاس یک عکس عوض شود، مسئله حمله به مدل عملا یک مسئله جستجو است. 

روشی که در مقاله بیان می‌شود به این صورت است که با توجه به این‌که عکس‌ها با مرز تصمیم فاصله کمی دارند، فرض می‌گیرد که جهت حرکت اهمیت کم‌تری داشته باشد و نیاز نیست بر روی این جهت خیلی تمرکز کرد. پس یک سری بردار جهت
($Q$)
 انتخاب شده و در هر گام یکی از این بردارها 
 $q \in Q$
 انتخاب شده و مقدار
$q \epsilon +‌ \delta$ 
به عکس اضافه می‌شود و بررسی می‌شود تا آیا امتیاز کلاس کاهش می‌یابد یا خیر. اگر امتیاز کاهش یافت مقدار 
$\delta$
با
$q \epsilon$
جمع زده می‌شود و اگر
 امتیاز کاهش نیافت عکس با بردار 
$ \delta -q \epsilon $
جمع زده می‌شود و مجدد چک می‌شود و درصورت کاهش امتیاز مقدار 
$\delta$
به روزرسانی می‌شود


برای این‌که تعداد کوئری‌ها کاهش یابد بهتر است از بردارهای متعامد استفاده شود تا همدیگر را خنثی نکنند. اگر بردارها متعامد نباشند، همچنان این حمله ممکن است ولی به تعداد کوئری‌های بیشتری نیاز خواهد داشت. 
\subsection{تفاوت \lr{SimBA} و \lr{SimBA-DCT}}
در روش 
\lr{SimBA}
 برای بردارهای کاندید، از مختصات دکارتی استفاده می‌شود. یعنی در هر دور آموزش رنگ یکی از پیکسل‌ها کم یا زیاد می‌شود. ولی در روش 
 \lr{SimBA-DCT}
 بردارها بجای تعریف روی فضای دکارتی روی فضای فرکانس تعریف می‌شوند. این عمل به این خاطر است که تحقیقات نشان داده که نویزها در فرکانس‌های پایین‌تر حالت خصمانه 
 \LTRfootnote{adversarial}
 بیشتری دارند. به این معنی که برای انسان قابلیت تشخیص کمتری دارند ولی بیشتر در خروجی مدل تاثیر دارند. پس یعنی با تغییر عکس در این حوزه می‌توان با فاصله کمتر، میزان امتیاز مدل را بیشتر کم کرد. 
\subsection{حملات امتیازمحور و حملات مبتنی بر برچسب }
در حملات مبتنی بر برچسب، مدل برای خروجی تنها گروهی که عکس در آن دسته‌بندی شده است را می‌دهد. ولی در حملات امتیاز محور، مدل احتمال تعلق عکس به هر دسته را می‌دهد. در حملات امتیاز محور از آنجایی که معیار پیوسته‌ای برای سنجش عملکرد حمله وجود دارد، تعداد کوئری‌ها کمتر است. مدل 
\lr{SimBA}
که در مقاله معرفی شده است یک مدل امتیاز محور است زیرا بر اساس این‌که با تغییر داده شده، امتیاز کلاس کم شده یا نه بردار خود را به‌روز می‌کند. 
\subsection{معیارهای ارزیابی}
در مقاله از سه معیار برای ارزیابی روش‌ها استفاده می‌شود:
\begin{itemize}
	\item 
	\textbf{نرخ موفقیت:}
	نرخ موفقیت بیان می‌کند در چند درصد مواقع حمله موفق شده که مدل را به اشتباه بیاندازد (یا در حالت حمله هدف‌دار، باعث شود تا مدل خروجی مطلوب را تولید کند.) هرچقدر این مقدار بیشتر باشد بهتر است زیرا حمله توانسته در موارد بیشتری مدل را به اشتباه بیندازد.  
	\item 
	\textbf{میانگین کوئری:}
	این مورد بیان می‌کند برای هر حمله، به طور میانگین چند کوئری استفاده شده است. منطقا هرچقدر این مقدار کمتر باشد بهتر است زیرا زمان و هزینه کمتری صرف می‌شود.
	\item 
	\textbf{میانگین $L_2$:}
	این مورد بیان می‌کند که عکس نهایی که باعث به اشتباه انداختن مدل می‌شود با عکس اولیه چقدر فاصله دارد. هرچقدر این عدد کمتر باشد یعنی مدل ضعیف‌تر بوده و حمله قوی تر بوده است. 
	

\end{itemize}
	\section{حملات جعبه سفید}
	\subsection{مدل و مجموعه داده }
	مدل و تصاویر با موفقیت بارگذاری شدند  و همانطور که در سوال آمده،‌دقت مدل روی داده‌های نمونه برابر با ۱۰۰ بود.
	\subsection{\lr{FGSM}}
	در روش 
	\lr{FGSM}
	ابتدا مقدار خطا برای ورودی و برچسب حساب می‌شود و سپس گرادیان خطا نسبت به ورودی حساب می‌شود سپس از این مقدار گرادیان فقط علامت آن در نظر گرفته می‌شود و در یک 
	$\epsilon$
	ضرب می‌شود و با عکس جمع می‌شود. اگر حمله هدف‌دار باشد، موقع محاسبه خطا، برچسب خروجی مدنظر داده می‌شود و موقع به روزرسانی عکس، در راستای نزول گرادیان حرکت می‌کنیم بجای صعود آن. 
	
	با انجام حمله 
	\lr{FGSM}
	روی نمونه‌ها، از ۲۵ عکس ورودی، تنها در ۲ مورد مدل همچنان دسته درست را پیش‌بینی می‌کرد و دقت مدل از ۱۰۰ به ۸ درصد افت کرد. 
	\begin{figure}[H]
		
		\centering
		\includegraphics[width=0.7\textwidth]{images/q2_fgsm_examples.png}
		\caption{
نمایش ۴ مورد از نمونه عکس اصلی، عکس دستکاری شده و نمایش بردار $\delta$ در 
\lr{FGSM}.
مقادیر دلتا از 
$-0.03$
تا 
$0.03$
بوده و برای آن‌که خروجی قابل تفسیری بدهد به صفر تا یک مقیاس شده است.
		}
		\label{q2:fgsm_samples}
	\end{figure}
	\subsection{\lr{PGD}}
	الگوریتم 
	\lr{PGD} 
	به نوعی، نسخه 
	\lr{iterative}
	الگوریتم 
	\lr{FGSM}
	است. به این صورت که در هر مرحله، مشابه 
	\lr{FGSM}
	مقدار گرادیان خروجی نسبت به ورودی حساب می‌شود. ولی بجای این‌که 
	$sign(gradient) * \epsilon$
	با عکس جمع شود و عکس نهایی خروجی داده شود، بلکه یک بردار اولیه 
	$\delta$
	با مقادیر صفر درنظر گرفته می‌شود و مقدار 
	$gradient * \alpha$
	به آن اضافه می‌شود. سپس چک می‌شود که این مقدار 
	$\delta$
	باعث می‌شود تا عکس دستکاری شده از همسایگی عکس اولیه خارج شود یا خیر و اگر خارج می‌شد 
	$\delta$
	کلیپ می‌شود و تاوقتی که مدل کلاس اشتباهی را انتخاب کند و یا تعداد دورها از ابر پارامتر بیشینه تعداد دور بیشتر شود، این کار تکرار می‌یابد. 
	
	با استفاده از این روش مدل تنها یک عکس را توانست درست پیش‌بینی کند و دقت آن به 
	$0.04$
	کاهش یافت. علیرغم این‌که این شیوه حمله از 
	\lr{FGSM}
 	قدرتمند تر است و توانست دقت را کمتر کند ولی حمله کندتری است و همانطور که از مقایسه شکل
	\ref{q2:fgsm_samples}
	و
	\ref{q2:pgd_samples}
	می‌توان یافت، بردار 
	$\delta$
	های یافت شده در دو الگوریتم بسیار شبیه به هم هستند. ولی برخلاف روش 
	\lr{FGSM}
	که در آن مقادیر هر کانال پیکسل‌ها باینری است (که در 
	$\epsilon$
	ضرب شده
	)
	در 
	\lr{PGD}
	این مقادیر می‌توانند مقادیری میانی بگیرند.
	
		\begin{figure}[H]
		
		\centering
		\includegraphics[width=0.7\textwidth]{images/q2_pgd_examples.png}
		\caption{
			نمایش ۴ مورد از نمونه عکس اصلی، عکس دستکاری شده و نمایش بردار $\delta$ در 
			\lr{PGD}.
			مقادیر دلتا از 
			$-0.03$
			تا 
			$0.03$
			بوده و برای آن‌که خروجی قابل تفسیری بدهد به صفر تا یک مقیاس شده است.
		}
		\label{q2:pgd_samples}
	\end{figure}

برای مقایسه بهتر دو حمله، میانگین فاصله 
$l_2$
در حملات نیز در جدول 
\ref{tab:l2_comp}
آمده است.
\begin{table}[H]
	\centering
	\caption{ میانگین فاصله $l_2$ در عکس‌های تولیدی}
	
	
	\begin{tabular}{r c}
		\toprule
		مدل &  میانگین فاصله $l_2$ \\
		\midrule
		
		\lr{FGSM} & $1.6628$ \\
		\lr{PGD} & $0.928$ \\
		\bottomrule
	\end{tabular}
	\label{tab:l2_comp}
\end{table}